% §17.3 The Trapdoor Computing Paradigm
%
% Monograph register: the paradigm (encrypt, compute on ciphertext, only
% key-holder decrypts), two-party server/client setup, confidentiality vs
% correctness on the perfect-filter line, Kronecker leakage over m
% queries, adaptive adversaries with a worked example.
% FIX: (1) notation collision in the worked example (the draft set A in
% {0,1}^n then wrote "n = |A|"); use N = |A| for dataset size throughout.
% (2) leakage-rate error (review M1): the draft claimed the adaptive
% adversary extracts O(log N) bits PER QUERY, which does not follow
% (seed-probing is O(1)/query; record-localization takes O(log N) queries
% TOTAL). Corrected: per-query extraction is O(1); the real gap is that a
% bounded budget of O(log N) crafted queries breaks the scheme outright,
% vs the iid attacker's O(fpr) bits/query that never recovers the seed.
% Cross-refs sec:kronecker, thm:kronecker, ax:element-indep,
% def:perfect-filter-line, sec:adversarial-inputs, def:cipher-map,
% sec:bernoulli-error-model, sec:fhe-sidebar, sec:maximizing-confidentiality.
% DEFINED: sec:trapdoor-paradigm.

\section{The Trapdoor Computing Paradigm}
\label{sec:trapdoor-paradigm}

The cipher map of \Cref{def:cipher-map} is the mathematical object; the
trapdoor computing paradigm is the operational setting in which it earns
its keep. The setting has three commitments: the data is encrypted before
it leaves the user, the computation runs entirely on the encrypted form,
and the result can be read only by a holder of a short secret key, the
trapdoor. The cipher map carries out the computation and the trapdoor
unlocks the answer; everything else, the plaintext records, the queries
that select among them, the intermediate results, stays hidden from any
party without the key.

\subsection*{The paradigm}

The paradigm predates the modern encryption-on-the-cloud framing.
Goldreich's work on software protection in the late 1980s asked how a
program could run on an untrusted machine without revealing what it
computed; the same question, in other language, drives oblivious RAM,
secure multiparty computation, and private information retrieval. What
unites them is the conviction that \emph{compute} and \emph{know} can be
separated: a machine can do useful work on data whose meaning it cannot
see. Trapdoor computing names the separation. The system has two
distinguished resources, an encrypted dataset and a trapdoor key, both
needed to read anything meaningful: the dataset without the key is opaque,
the key without the dataset useless. Computation is the operation of
cipher maps on the encrypted dataset, producing encrypted outputs the key
holder can decrypt and no one else can. The Bernoulli framework enters at
decryption: because the cipher map is approximate, the decrypted output
is a Bernoulli observation rather than a deterministic fact, and the
chapter 4 axioms govern how it behaves across queries. The cryptographic
side sets up the problem (compute without revealing); the Bernoulli side
tracks the consequence (read with controlled error).

\subsection*{Two-party setup}

The cleanest instance has a \emph{server} holding the encrypted dataset
and running the cipher maps, and a \emph{client} holding the trapdoor and
issuing queries. The server sees ciphertexts in and ciphertexts out,
nothing else. The encryption keeps the query private; the cipher map's
bounded error keeps the dataset partially hidden even from a client who
probes many times. The standard threat model is the
\emph{honest-but-curious} server, which follows the protocol but records
the encrypted traffic for later analysis; bounding what that analysis can
recover is the Bernoulli framework's job. A malicious server that deviates
from the protocol needs cryptographic integrity checks beyond this
chapter. The client's recovery quality is set by the $(\fprate, \fnrate)$
profile: an $\fnrate = 0$ map loses no real positives but admits some
false ones, an $\fprate = 0$ map trusts every positive but misses some
real ones.

\subsection*{Confidentiality versus correctness}

The trade-off that distinguishes the paradigm from cleartext computation
is between confidentiality and correctness. Tight error rates make the
client's view faithful but turn each near-deterministic answer into a
side channel: a server watching many such queries ties encrypted queries
to the records they match and reconstructs the plaintext distribution.
Loose rates make each query noisier for the client but push the
encrypted-output distribution toward uniform across the dataset, defeating
frequency analysis. The $(\fprate, \fnrate)$ point is therefore a
confidentiality knob as well as a correctness knob, and the
perfect-filter line of \Cref{def:perfect-filter-line} reads naturally in
this light: an $\fnrate = 0$ map concentrates its confidentiality cost in
the single parameter $\fprate$ (every real positive returned, recall
preserved), an $\fprate = 0$ map in $\fnrate$ (every positive answer
correct, but real positives missed). The first suits recall-sensitive
retrieval, the second precision-sensitive contexts.
\Cref{sec:maximizing-confidentiality} returns with a quantitative measure.

\subsection*{Multiple queries: leakage accumulates}

A single query reveals little; a million can reveal much. Under
\Cref{ax:element-indep} (Axiom 1) the error events of $m$ independent
queries are mutually independent, and the joint distribution factors as
the Kronecker product of $m$ copies of the per-query channel matrix
(\Cref{thm:kronecker}). The same algebra describes the trapdoor stream,
and any single-query confidentiality bound lifts through tensor algebra to
the stream. The independence cuts both ways: the client's recovery
sharpens as $m^{1/2}$ by standard concentration, and the server's
information about the plaintext grows linearly in $m$ for any map with
positive single-query leakage. The two scalings are the same
data-processing inequality read from opposite sides, the Kronecker
structure making the duality exact. A system with a finite query budget
can be analyzed by applying the per-query bounds to the $m$-fold Kronecker
product; an unbounded budget needs a different approach, either rotating
the trapdoor key to reset accumulated leakage or designing for the
short-horizon regime where asymptotic confidentiality has not yet decayed
to zero.

\subsection*{Adaptive adversaries}

The independence the Kronecker analysis depends on fails when the stream
is adaptive: the adversary chooses each query from the encrypted outputs
seen so far. This is exactly the adversarial-input setting of
\Cref{sec:adversarial-inputs}, and the trapdoor paradigm inherits its
failure modes, with the encrypted outputs (not plaintext ones) driving
the strategy. The bounds shift from average-case (iid Kronecker products)
to worst-case: a map at $\fprate = 0.01$ on iid queries may reach
$\fprate$ near $1$ on a query crafted from the preceding outputs. The
response, as in chapter 16, is worst-case accounting: bound the maximum
plaintext bits any adversary extracts per query and budget around that.

Consider a trapdoor membership map on a plaintext set $A$ of $N = |A|$
records, with $\fprate = 2^{-k}$ and $\fnrate = 0$, realized by a
Bloom-style construction with $k$ keyed hash functions whose seed is the
trapdoor. Under iid queries the behavior is benign: each random query has
independent probability $2^{-k}$ of a false positive, and after $m$
queries the server has gained $O(m)$ bits about $A$, spread uniformly
enough that no record is pinned down. Under an adaptive adversary the
picture changes, not in the per-query rate but in what a bounded budget
of crafted queries buys. The $\fprate = 2^{-k}$ comes from $k$ hash bits
matching at once; an adversary crafts queries that probe the bits
individually, each extracting about one bit of the seed, so after $O(k)$
such queries the seed is recovered and the entire dataset decrypts, and a
binary search over the records localizes a chosen record's membership in
$O(\log N)$ queries. The per-query extraction stays $O(1)$ bits
throughout; the danger is that this bounded budget, logarithmic in the
dataset size, breaks the scheme outright, where the iid attacker gains
only $O(\fprate)$ bits per query and never recovers the seed. The gap
between total compromise in a logarithmic number of crafted queries and
an $O(\fprate)$-bit-per-query trickle is the gap between worst-case and
average-case, and the fix, as chapter 16 prescribed, is a cipher map no
bounded budget of crafted queries can break. Cryptographic hash families
(pseudorandom functions) achieve this: a pseudorandom seed cannot be
recovered by probing, so the worst-case attacker gains only one bit of
best-effort guessing per query, matching the average case to a constant;
the FHE sidebar that follows gives an even stronger guarantee, worst-case
leakage exactly zero. Both
the worst-case story (chapter 16) and the average-case story (chapters 4
through 14) are needed in the trapdoor setting: the average case for
sizing under benign workloads, the worst case for security against
adversarial ones.
